\documentclass[11pt,a4paper]{article}
\usepackage[margin=23mm]{geometry}
\usepackage{amsmath,amssymb,booktabs,array}
\input{glyphtounicode}
\pdfgentounicode=1
\usepackage[hidelinks]{hyperref}
\hypersetup{hypertexnames=false}
\setlength{\parindent}{0pt}
\setlength{\parskip}{5pt}
\newcommand{\pos}[1]{[#1]_+}
\renewcommand{\labelitemi}{\(\bullet\)}
\begin{document}
\begin{center}
{\Large Selected W-F: joined forward qualification}\\[4pt]
16 October 2026 --- the original constant-advice domain; no traffic experiment
\end{center}

\textbf{Status.} The user selects W-F and closes J4-W1: merging work has
no direct advice dependence at fixed represented physical state. The
conditional whole-model forward argument below is delivered at
\textbf{WF-Q1}, an evidence checkpoint. Physical MP6 and the dependent
action, gradient and reference contracts remain open. Earlier proofs keep
their stated scope; this is a new deduction for the selected joint member.

\section*{1. Primary-domain assumptions and selected expression}
Retain the accepted H1 preparation and initially empty finite T1 terminal:
fixed $L=L_p+b$, initial footprint $[0,\ell]$, $0<\ell\leq L_p$,
initial densities $e_0,a_0=b_0\geq0$, $r=e_0+a_0$, and
$M=\ell(e_0+2a_0)>0$. There are no postdecision external arrivals.
One constant common action lies in the original interval $[v_-,v_+]$.
In addition to H1's admitted initial and handoff conditions, retain
\begin{equation}
 v_+r<R\leq C(v_-),\quad v_+b_0\leq R_B\leq C(v_-),
 \quad M<k_jb,\quad w,k_j,h>0.\tag{1}
\end{equation}
With $K=wk_j$, define
\begin{equation}
 C(v)=\frac{Kv}{w+v},\quad U(f)=\frac{wf}{K-f},\quad
 P_j(t)=F_j(t-b/v),\quad q_j=P_j-Y_j,\quad N_j=F_j-Y_j.\tag{2}
\end{equation}
The far lane retains $Y_B=P_B$. For A/E, $g_j=\dot Y_j$ and
$\lambda_j=\dot P_j$. Write $B(t)$ for shared service, distinct from
the far-lane index. Free marks are $u_j=v$; queued marks are $U(g_j)$.

Set fixed $v_\star=v_+$ and $C_\star=C(v_\star)$ once, and use
\begin{equation}
 c_A=1,\quad c_E=c_0+\chi g_A\pos{u_A-u_E},\quad
 c_0=1+\delta\geq1,\quad \chi=\frac{\beta}{C_\star v_\star},\quad
 \eta=c_0-\beta(1+v_+/w)>0.\tag{3}
\end{equation}
$\beta\geq0$ and all physical coefficients are fixed. $\chi$ has units
s$^2$/(veh-eq m); $c_j$ is dimensionless work per unit mass. The anchor
is a fixed coefficient convention, not a variable to recompute with advice
or an optimizer's search interval. At fixed $g_j,u_j$, (3) has zero direct
advice derivative. The prior conservative root bound and
$c_0\leq c_E\leq c_{\max}:=c_0+\beta$ remain sufficient.

\newpage
\section*{2. The retained causal node and recovery account}
For a queued E, let $\mathcal H(g;a,u)$ denote
\begin{equation}
 \mathcal H(g;a,u)=g\{c_0+\chi a\pos{u-U(g)}\},\qquad
 \partial_g\mathcal H\geq\eta>0.\tag{4}
\end{equation}
Its unique root for available work $x$ lies between
$x/c_{\max}$ and $x/c_0$, and below $C(v)$. The selected S-E-P
branches, on their feasible domains, are:
\begin{center}
\begin{tabular}{>{\raggedright\arraybackslash}p{28mm}>{\raggedright\arraybackslash}p{115mm}}
\toprule
Current conditions & Actual service\\
\midrule
A and E free & $g_A=\lambda_A$, $g_E=\lambda_E$ when
  $\lambda_A+c_0\lambda_E\leq B$.\\[3pt]
A free, E queued & $a=\lambda_A\leq B$, $g_A=a$ and
  $\mathcal H(g_E;a,v)=B-a$.\\[3pt]
A and E queued & $g_A=B/2$ and
  $\mathcal H(g_E;B/2,U(B/2))=B/2$.\\[3pt]
A queued, E free & $g_E=\lambda_E$, $g_A=B-c_0\lambda_E$ when
  $c_0\lambda_E\leq B/2$; otherwise E's offer also has a shortfall.\\
\bottomrule
\end{tabular}
\end{center}
These are the existing equal-work/unused-service template and free-A
priority, not a new allocation. A positive queue follows
$\dot q_j=\lambda_j-g_j$; at zero, only available mass can pass.
An empty stream stays free while that continuation is feasible. A strict
loss of free feasibility initiates its queued branch. At a zero hit, test
free admission again. Equality alone is not a queue-birth cause. This is
the previously selected causal guard, stated as first loss of feasibility
for the proof; no queue floor, latch or numerical selection is introduced.

R-E and O-R remain unchanged. With positive acceleration $a_r$ and fixed
$m=m_{\rm ref}=1$ vehicle-equivalent, use
\begin{align}
 c(u)&=\frac{Ku}{w+u},\quad u_R=U(R),\quad
 D(\theta;u)=\pos{R-c(\sqrt{(w+u)^2+2wa_r\theta}-w)},\nonumber\\
 p(t)&=\frac1m\int_{[Y_E(t)-m]_+}^{Y_E(t)}
       D(t-T(\xi);u_E(T(\xi)))d\xi,\quad B(t)=\pos{R-p(t)}.\tag{5}
\end{align}
$T$ is the admitted E mass-rank inverse; zero-rate plateaus have zero
admission measure. New mass adds profiles and retires oldest weight after
the window fills; all retained profiles age physically. Only E deposits
new loss. All marked history is predicted local state. Releasing a queue
does not reset that history, and retiring history does not remove traffic.

\newpage
\section*{3. The H1 pulse limits E to one waiting episode}
Eligible arrivals start at $t_f=(L-\ell)/v$ and end at $t_r=L/v$.
Inside this interval H1 implies
\begin{equation}
 \lambda_A+\lambda_E=s_v:=vr<R,\quad
 \lambda'_A=-\lambda'_E\in[0,L_A],\quad
 L_A:=v_+^2e_0/h.\tag{6}
\end{equation}
Indeed $\lambda_E=v e_{H1}(t-b/v)$ and
$\dot e_{H1}=-f/h$, $0\leq f\leq v e_{H1}\leq ve_0$.
There are only the front and rear jumps; H1 branch switches preserve
continuous rates. No new initial/demand restriction is used in (6).

Initially $B=R$ and A's offer is below $R$. If
\begin{equation}
 W_f(v):=\lambda_A(t_f+)+c_0\lambda_E(t_f+)\leq R,\tag{7}
\end{equation}
all free traffic passes. Since $v\geq u_R$, free E has $D=0$.
Then $B=R$ and $\lambda_A+c_0\lambda_E=s_v+\delta\lambda_E$
is nonincreasing; the free solution persists through the cohort. If (7)
fails strictly, A still fits, and E's queued root satisfies
$g_E\leq(R-\lambda_A)/c_0<\lambda_E$. E therefore begins waiting
at the front with a compatible positive right-hand derivative.

Suppose that positive E queue later first reaches zero, during the pulse.
If A is free, its pre-clearance root $g_E^Q$ gives
\begin{equation}
 \lambda_E\leq g_E^Q\quad\Longrightarrow\quad
 \lambda_A+c_0\lambda_E\leq B.\tag{8}
\end{equation}
If A remains queued, $\mathcal H(g_E^Q;B/2,U(B/2))=B/2$ gives
\begin{equation}
 c_0\lambda_E\leq c_0g_E^Q\leq B/2.\tag{9}
\end{equation}
Thus free E is feasible in either case. New free E profiles are zero.
In (5), aging and retirement can only reduce existing loss; hence $B$
is nondecreasing as long as E stays free. Since $\lambda_E$ decreases,
(9) remains valid while A queues. If A clears, its queued service before
clearance was $B-c_0\lambda_E$, so free joint admission satisfies (8).
Thereafter the joint offered work decreases while $B$ increases. This
continuation is self-consistent and cannot activate E again.

Simultaneous A/E clearance also gives joint free feasibility: before the
hit $\lambda_A\leq B/2$ and $c_0\lambda_E\leq B/2$. After the rear
there are no arrivals, so neither queue can re-form. Consequently E has
at most one waiting episode: it is absent, or begins at the front and
eventually clears. This is specific to the retained monotone H1 pulse,
neutral initial account and R-E; arbitrary traffic is not covered.

\newpage
\section*{4. Positive budget and finite A switching}
From (3), admitted E rate is at most $G=R/c_0$. The unchanged profile
is nonincreasing in age and mark, bounded by $R$, and strictly less than
$R$ at positive age when $a_r>0$. The earlier rearrangement argument gives
\begin{equation}
 p\leq p_{\max}:=\frac Gm\int_0^{m/G}D(\theta;0)d\theta<R,
 \qquad B\geq B_{\min}:=R-p_{\max}>0.\tag{10}
\end{equation}
This is derived, not an added floor. $B_{\min}$ is common to all actions
for fixed physical parameters. Its value can be very small.

The injection, aging and retirement update in the earlier O-R note yields
an a.e. Lipschitz bound
\begin{equation}
 |\dot B|\leq L_B:=2GR/m+a_rk_j.\tag{11}
\end{equation}
The retained mass is at most $m$, $|D_\theta|\leq a_rk_j$, and each
boundary term is at most $GR/m$. No inverse-rate lower bound is needed.
This bound permits speed-mark jumps at queue events.

Before E clears, any A birth inside the pulse occurs when the free offer
first loses feasibility, at $\lambda_A=B$. Switching to queued cooperation
gives $\dot q_A=B/2\geq B_{\min}/2>0$. If A subsequently clears
while E still queues, continuity of the pre-clearance derivative gives
$\lambda_A\leq B/2$. For $y=B-2\lambda_A$, its value at birth is
$-B\leq-B_{\min}$ and its value at clearance is nonnegative. By
(6),(11), a birth-to-clearance interval is at least
\begin{equation}
 \frac{B_{\min}}{L_B+2L_A}.\tag{12}
\end{equation}
At A clearance, the free margin $x=B-\lambda_A$ is at least
$B_{\min}/2$. The next A birth needs $x=0$; thus a clearance-to-birth
interval is at least $B_{\min}/[2(L_B+L_A)]$. A common smaller bound is
\begin{equation}
 \tau_0=\frac{B_{\min}}{2(L_B+2L_A)}>0.\tag{13}
\end{equation}
The bound applies to consecutive A events while E queues and the arrival
pulse continues. E clearance and the one external pulse end add finitely
many exceptions. Once E is free, the previous section excludes further E
births, and a cleared A cannot re-form because joint free work fits and
decreases. After the rear, positive queues only drain.

Therefore A events cannot accumulate in finite time. Tangential zero hits
use the already selected free-admission test; the post-clearance margin
still holds. The positive separation is a proof consequence, not an
implemented dwell timer or additional historical allocation preference.

\newpage
\section*{5. Global continuation, storage and complete clearance}
\textbf{Causal construction.} On each valid branch, (4) and the bounded
derivative of $U$ on $[0,R]$ make flows and speed marks locally Lipschitz
in $B$ and continuous arrival data. The earlier moving-cutoff O-R estimate
for histories with the same past is
\begin{equation}
 |\Delta p(t)|\leq\frac1m\left(
 5R\int|\Delta g_E|ds+k_jG\int|\Delta u_E|ds\right).\tag{14}
\end{equation}
This gives short-interval Volterra contraction on a valid branch, including
cutoff movement and old mark jumps. Continue to its next guard. At the
front E shortfall, queue birth has the strict compatible sign shown above.
At A birth, the new right-hand derivative is positive. A zero hit admits
free A; E zero hits admit free E by (8)--(9). Simultaneous clearance admits
both. Equality contacts do not themselves cause birth; keep the maximal
feasible free continuation, as the selected guard requires. Point values
of flux at a switch do not alter continuous cumulative mass.

Boundedness prevents finite-time blow-up within branches. Before the rear,
E switches at most twice and repeated A events are separated by (13), so
branch continuation cannot require infinitely many switches on that finite
pulse. After the rear, no new queues can be born; there are at most two
clearances. This proves a unique causal, piecewise-branch continuation for
every original advice under (1)--(3) and the selected causal convention.
No transversality assumption or advice subinterval is needed for forward
existence; smooth dependence on advice is a separate issue.

If $Q=q_A+q_E>0$ after $t_r$, the node uses the available work and
$g_A+g_E\geq B/c_{\max}\geq B_{\min}/c_{\max}$. Therefore
\begin{equation}
 T_{\rm clear}(v)\leq\frac Lv+
        \frac{c_{\max}Q(L/v)}{B_{\min}}
 \leq T_{\max}:=\frac L{v_-}+\frac{c_{\max}M}{B_{\min}}<\infty.\tag{15}
\end{equation}
Thus the entire cohort clears. Afterwards profiles simply age out within
the unchanged finite support. The bound does not assert practical runtime.

H1 entrance flows fit $C(v)$ and $M<k_jb$. The retained T1 vacancy proof
therefore protects the handoff for these admissible departures. Physical
stock remains $N_j=F_j-Y_j$, with $0\leq Y_j\leq P_j\leq F_j$.
If $s_{\rm occ}$ is upstream occupied length and $\rho_0=e_0+2a_0$,
\begin{equation}
 s_{\rm occ}\rho_0+\sum_jN_j+\sum_jY_j=M,\qquad
 J(v)=\frac{M(L-\ell/2)}v+\int_0^\infty(q_A+q_E)dt.\tag{16}
\end{equation}
No work weight changes traffic mass, and post-exit recovery is not local
vehicle-time. The result is not a general invariant LTM network theorem.

\newpage
\section*{6. Information, negative controls and sensitivity limits}
Known $M,\ell,E_0$ and continuing symmetry reconstruct
\begin{equation}
 e_0=E_0/\ell,\qquad a_0=b_0=(M-E_0)/(2\ell).\tag{17}
\end{equation}
Together with known empty terminal/account and fixed parameters, uniqueness
now gives identical full forward prediction for A+E and explicit lane G
for every original advice. Their predicted history is not an observation.
This does not imply that the current meter alone is sufficient dynamic
state. The earlier pre-decision collision stays upstream of an empty
terminal and is unchanged; no old action/gradient certificate is transferred.

\textbf{A sharper negative control.} If $\delta=0$, (7) always holds
since $\lambda_A+\lambda_E=s_v<R$. Hence no terminal queues or recovery
are activated, regardless of $\beta$: free A and E have equal marks.
The entire primary objective is $M(L-\ell/2)/v$, so fastest advice is
common to all preparations. No E is another such control. Immediate
maneuver work is therefore the trigger in this particular underloaded
model, not a universally necessary congestion mechanism.

\textbf{Front threshold.} With fixed $z_f=L_p-\ell$, H1 comparison gives
$e_f(v)=e_{H1}(z_f;v)$ nondecreasing in $v$: $k_c(v)$ decreases, so
the receiving-limited reduction becomes no stronger. Thus
\begin{equation}
 W_f(v)=v\{r+\delta e_f(v)\}
 \text{ is strictly increasing for }r>0.\tag{18}
\end{equation}
There is at most one transition between a free and initially queued cohort
as advice increases. It may lie outside the action interval. Equality
belongs to free passage under the selected guard.

Forward uniqueness does not imply differentiability at that threshold.
If it is interior, with $a=\lambda_A>0$, $e=\lambda_E>0$ and
$a+c_0e=R$, then for $\chi>0$,
\begin{equation}
 \mathcal H(e;a,v)-(R-a)=\chi a e\{v-U(e)\}>0.\tag{19}
\end{equation}
The queued-root limit is strictly less than the free flow $e$. Speed marks
and later recovery can therefore change abruptly across activation.
Equation (19) is a local flux statement, not a proof of a jump in complete
$J$. An ordinary smooth derivative across it must not be presumed. Stable
event sensitivities, grazing/boundary cases, cutoff motion and full-objective
tests remain in the next evidence work; no AD/FD run is performed here.

\textbf{Scope/checkpoint.} The planned forward block is complete at WF-Q1.
The older regular A-delay/active-renewal witness is exactly preserved at the
fixed reference advice. Physical speed identification, work/recovery
partition, joint scales/applied MP6 and independent usefulness remain open.
Carey's work interpretation, Leclercq's isolated recovery context and Jin's
separate discontinuous-boundary model motivate ingredients, not this proof.
Step 3 is open; Step 4 and implementation/experiments remain unopened.
\end{document}
